% ===========================================================================
%  Chapitre 11 — Suites arithmétiques et suites géométriques
%  RAPE, quatrième période — hors tableau des contenus du PE 2026
% ===========================================================================
\bmtchapitre{Suites arithmétiques et suites géométriques}
  {Reconnaître une suite arithmétique ou géométrique, en donner la raison, le
   terme général, la somme de $k$ termes consécutifs et la limite ; utiliser
   une suite auxiliaire pour ramener une suite récurrente à l'une de ces deux
   familles.}
  {Suites numériques \textbullet\ environ 6 heures}

Voici le chapitre le plus rentable du livre après celui du logarithme. Le
premier exercice du baccalauréat porte sur les suites à chaque session
examinée pour cet ouvrage — de 1999 à 2026, sans exception — et il vaut cinq
points, soit un quart de la note. Mieux : sa structure ne varie pratiquement
pas. Une suite récurrente est donnée ; on pose une suite auxiliaire, presque
toujours notée $V_{n}$ ; on démontre qu'elle est géométrique ; on en déduit
$V_{n}$, puis $U_{n}$, puis une somme, puis une limite.

Autrement dit : cinq points par an, pour une seule méthode apprise une fois.

% ---------------------------------------------------------------------------
\begin{revision}
Le chapitre 10, et les propriétés du logarithme.

\begin{enumerate}[label=\textbf{\arabic*.}, leftmargin=*, itemsep=0.35em]
  \item $u_{0} = 7$ et $u_{n+1} = u_{n}+3$ : calculer $u_{1}$, $u_{2}$ et
        $u_{3}$.
  \item $v_{0} = 2$ et $v_{n+1} = 5v_{n}$ : calculer $v_{1}$ et $v_{2}$.
  \item Que vaut $\limn q^{n}$ pour $q = 0{,}4$ ? pour $q = 1{,}4$ ?
  \item Calculer $1+2+3+\cdots+10$.
  \item Simplifier $\ln\left(a^{n}\right)$ pour $a>0$.
  \item Écrire $e^{3n-1}$ sous la forme $a\times q^{n}$.
\end{enumerate}
\end{revision}

\begin{automatismes}
Sans calculatrice.

\begin{enumerate}[label=\textbf{\alph*)}, leftmargin=*, itemsep=0.25em]
  \item Suite arithmétique, $u_{0} = 1$, $r = 4$ : donner $u_{3}$
  \item Suite géométrique, $u_{0} = 3$, $q = 2$ : donner $u_{3}$
  \item Terme général d'une suite arithmétique de premier terme $u_{0}$
  \item Terme général d'une suite géométrique de premier terme $u_{0}$
  \item Nombre de termes de $u_{0}+u_{1}+\cdots+u_{n}$
  \item Nombre de termes de $u_{5}+u_{6}+\cdots+u_{20}$
  \item $1+q+q^{2}+\cdots+q^{n}$ pour $q \neq 1$
  \item $\limn \left(\frac{1}{4}\right)^{n}$
\end{enumerate}
\end{automatismes}

\begin{corrige}[automatismes]
a) $13$ \quad b) $24$ \quad c) $u_{n} = u_{0}+nr$ \quad
d) $u_{n} = u_{0}q^{n}$ \quad e) $n+1$ \quad f) $16$ \quad
g) $\frac{1-q^{n+1}}{1-q}$ \quad h) $0$.
\end{corrige}

% ===========================================================================
\section{Les suites arithmétiques}

\begin{definition}[suite arithmétique]
La suite $\left(u_{n}\right)$ est \textbf{arithmétique} s'il existe un réel
$r$, appelé \textbf{raison}, tel que
\[
  u_{n+1} = u_{n}+r \qquad \text{pour tout entier } n .
\]
On passe d'un terme au suivant en \emph{ajoutant} toujours le même nombre.
\end{definition}

\begin{propriete}[terme général, variation, limite]
Pour une suite arithmétique de raison $r$ :
\[
  u_{n} = u_{0}+nr, \qquad\text{et plus généralement}\qquad
  u_{n} = u_{p}+(n-p)\,r .
\]
La suite est strictement croissante si $r>0$, strictement décroissante si
$r<0$, constante si $r = 0$. Sa limite est $+\infty$ si $r>0$, $-\infty$ si
$r<0$, et $u_{0}$ si $r = 0$.
\end{propriete}

\begin{theoreme}[somme de termes consécutifs]
\[
  S = \left(\text{nombre de termes}\right)
      \times \frac{\text{premier terme}+\text{dernier terme}}{2} .
\]
En particulier $u_{0}+u_{1}+\cdots+u_{n} = (n+1)\dfrac{u_{0}+u_{n}}{2}$.
\end{theoreme}

\begin{avertissement}
Le nombre de termes de $u_{0}+\cdots+u_{n}$ est $n+1$, et non $n$ : on
compte le rang $0$. Plus généralement, de $u_{p}$ à $u_{n}$ il y a
$n-p+1$ termes. Cette erreur de comptage est la plus fréquente de tout
l'exercice 1.
\end{avertissement}

\begin{methode}{montrer qu'une suite est arithmétique}
On calcule $u_{n+1}-u_{n}$ et l'on montre que le résultat est un
\textbf{nombre constant}, indépendant de $n$. On conclut en écrivant : « la
suite est arithmétique de raison $r = \dots$ et de premier terme
$u_{0} = \dots$ ». Les deux renseignements sont exigés.
\end{methode}

\begin{exemple}[une suite arithmétique cachée]
Soit $\left(U_{n}\right)$ définie par $U_{0} = 2$ et
$2U_{n+1}+1 = U_{n+1}+U_{n}$.

En retranchant $U_{n+1}$ des deux côtés : $U_{n+1}+1 = U_{n}$, donc
$U_{n+1} = U_{n}-1$. La différence $U_{n+1}-U_{n}$ vaut $-1$ : la suite est
arithmétique de raison $r = -1$ et de premier terme $U_{0} = 2$. Son terme
général est $U_{n} = 2-n$.
\end{exemple}

\begin{entrainement}[suites arithmétiques]
\exo[\niveau{1}] $u_{0} = -3$, $r = 5$ : donner $u_{n}$ puis $u_{12}$.

\exo[\niveau{1}] $u_{n} = 7-4n$ : montrer que la suite est arithmétique et
donner sa raison.

\exo[\niveau{2}] Suite arithmétique telle que $u_{3} = 11$ et $u_{8} = 31$ :
déterminer $r$ et $u_{0}$.

\exo[\niveau{2}] $u_{0} = 1$ et $r = -2$ : déterminer $n$ tel que
$u_{n} = -99$.

\exo[\niveau{2}] Calculer $S = u_{0}+u_{1}+\cdots+u_{50}$ pour la suite de
l'exercice précédent.

\exo[\niveau{3}] Suite arithmétique telle que $u_{34} = u_{15}+57$ :
déterminer la raison.

\exo[\niveau{3}] Pour cette même suite, sachant que $u_{22} = 65$, calculer
$u_{15}$, $u_{34}$, puis $S = u_{15}+u_{16}+\cdots+u_{34}$.
\end{entrainement}

\begin{corrige}[suites arithmétiques]
\textbf{1.} $u_{n} = -3+5n$ ; $u_{12} = 57$.
\textbf{2.} $u_{n+1}-u_{n} = -4$ : arithmétique de raison $-4$.
\textbf{3.} $u_{8} = u_{3}+5r$ donne $31 = 11+5r$, donc $r = 4$ ; puis
$u_{0} = u_{3}-3r = 11-12 = -1$.
\textbf{4.} $1-2n = -99$ donne $n = 50$.
\textbf{5.} $51$ termes, premier $1$, dernier $-99$ :
$S = 51\times\frac{1-99}{2} = 51\times(-49) = -2\,499$.
\textbf{6.} $u_{34} = u_{15}+19r$, donc $19r = 57$ et $r = 3$.
\textbf{7.} $u_{15} = u_{22}-7r = 65-21 = 44$ et
$u_{34} = u_{22}+12r = 65+36 = 101$. La somme compte $20$ termes :
$S = 20\times\frac{44+101}{2} = 20\times72{,}5 = 1\,450$.
\end{corrige}

% ===========================================================================
\section{Les suites géométriques}

\begin{definition}[suite géométrique]
La suite $\left(u_{n}\right)$ est \textbf{géométrique} s'il existe un réel
$q \neq 0$, appelé \textbf{raison}, tel que
\[
  u_{n+1} = q\,u_{n} \qquad \text{pour tout entier } n .
\]
On passe d'un terme au suivant en \emph{multipliant} toujours par le même
nombre.
\end{definition}

\begin{propriete}[terme général, somme, limite]
Pour une suite géométrique de raison $q$ :
\[
  u_{n} = u_{0}\,q^{n}, \qquad
  u_{0}+u_{1}+\cdots+u_{n} = u_{0}\,\frac{1-q^{n+1}}{1-q}
  \quad (q \neq 1).
\]
Sa limite se lit sur $q$ : $0$ si $-1<q<1$ ; $u_{0}$ si $q = 1$ ;
$\pm\infty$ selon le signe de $u_{0}$ si $q>1$ ; pas de limite si
$q \leq -1$.
\end{propriete}

\begin{methode}{montrer qu'une suite est géométrique}
On calcule le \textbf{quotient} $\dfrac{u_{n+1}}{u_{n}}$ — ou, plus sûr
encore, on exprime $u_{n+1}$ en fonction de $u_{n}$ et l'on fait apparaître
un facteur constant. On conclut en donnant la raison \emph{et} le premier
terme.
\end{methode}

\begin{exemple}[une géométrique écrite avec une exponentielle]
Soit $V_{n} = e^{2-n}$. Alors
\[
  \frac{V_{n+1}}{V_{n}} = \frac{e^{1-n}}{e^{2-n}} = e^{-1},
\]
un nombre constant : la suite est géométrique de raison $q = e^{-1}$ et de
premier terme $V_{0} = e^{2}$. Comme $0<e^{-1}<1$, elle converge vers $0$.
\end{exemple}

\begin{remarque}
C'est exactement la situation que posent les sujets malgaches : on y trouve
$V_{n} = e^{2-n}$, $V_{n} = e^{3n-1}$, $V_{n} = e^{3-2n}$ ou encore
$V_{n} = e^{-U_{n}}$. Dans tous les cas, le même geste : écrire le quotient
de deux termes consécutifs et simplifier les exponentielles.
\end{remarque}

\begin{entrainement}[suites géométriques]
\exo[\niveau{1}] $u_{0} = 5$, $q = 2$ : donner $u_{n}$ et $u_{6}$.

\exo[\niveau{1}] $u_{n} = 3\left(\frac58\right)^{n}$ : exprimer $u_{n+1}$ en
fonction de $u_{n}$ et donner la limite.

\exo[\niveau{2}] $V_{n} = e^{3n-1}$ : montrer que la suite est géométrique,
préciser sa raison et son premier terme.

\exo[\niveau{2}] Suite géométrique à termes positifs, de raison
$q = \frac34$, telle que $V_{3} = \frac{27}{8}$ : justifier que $V_{0} = 8$.

\exo[\niveau{2}] Calculer $S = 1+\frac12+\frac14+\cdots+\frac{1}{2^{10}}$.

\exo[\niveau{3}] $u_{n} = 2\left(\frac13\right)^{n}$ : calculer
$S_{n} = u_{0}+\cdots+u_{n}$ en fonction de $n$, puis $\limn S_{n}$.
\end{entrainement}

\begin{corrige}[suites géométriques]
\textbf{1.} $u_{n} = 5\times2^{n}$ ; $u_{6} = 320$.
\textbf{2.} $u_{n+1} = \frac58u_{n}$ ; la limite vaut $0$.
\textbf{3.} $\frac{V_{n+1}}{V_{n}} = e^{3}$ : géométrique de raison
$e^{3}$, de premier terme $V_{0} = e^{-1}$.
\textbf{4.} $V_{3} = V_{0}q^{3} = V_{0}\times\frac{27}{64}$, donc
$V_{0} = \frac{27}{8}\times\frac{64}{27} = 8$.
\textbf{5.} $11$ termes, premier $1$, raison $\frac12$ :
$S = \frac{1-\left(\frac12\right)^{11}}{1-\frac12}
= 2\left(1-\frac{1}{2\,048}\right) = \frac{2\,047}{1\,024}$.
\textbf{6.} $S_{n} = 2\times\frac{1-\left(\frac13\right)^{n+1}}{\frac23}
= 3\left(1-\left(\frac13\right)^{n+1}\right)$, de limite $3$.
\end{corrige}

\begin{qcm}
\exo Une suite telle que $u_{n+1} = u_{n}-7$ est :
\textbf{a)} géométrique de raison $-7$ \quad
\textbf{b)} arithmétique de raison $-7$ \quad
\textbf{c)} constante \quad \textbf{d)} ni l'une ni l'autre \reponseqcm{b}

\exo Pour une suite géométrique de raison $q = -3$, la limite est :
\textbf{a)} $0$ \quad \textbf{b)} $-\infty$ \quad \textbf{c)} $+\infty$
\quad \textbf{d)} il n'y en a pas \reponseqcm{d}

\exo La somme $u_{3}+u_{4}+\cdots+u_{20}$ comporte :
\textbf{a)} $17$ termes \quad \textbf{b)} $18$ termes \quad
\textbf{c)} $20$ termes \quad \textbf{d)} $23$ termes \reponseqcm{b}

\exo Si $V_{n} = U_{n}-4$ est géométrique de raison $\frac13$ et
$V_{0} = 2$, alors $U_{n}$ vaut :
\textbf{a)} $2\left(\frac13\right)^{n}$ \quad
\textbf{b)} $2\left(\frac13\right)^{n}+4$ \quad
\textbf{c)} $2\left(\frac13\right)^{n}-4$ \quad
\textbf{d)} $4\left(\frac13\right)^{n}$ \reponseqcm{b}

\exo $W_{n} = \ln\left(U_{n}\right)$ où $\left(U_{n}\right)$ est géométrique
de raison $q>0$. Alors $\left(W_{n}\right)$ est :
\textbf{a)} géométrique de raison $\ln q$ \quad
\textbf{b)} arithmétique de raison $\ln q$ \quad
\textbf{c)} constante \quad \textbf{d)} sans nature particulière
\reponseqcm{b}
\end{qcm}

% ===========================================================================
\section{La suite auxiliaire : la méthode de l'exercice 1}

\begin{methode}{le schéma que le baccalauréat répète chaque année}
Une suite $\left(U_{n}\right)$ est donnée par $U_{0}$ et une relation
$U_{n+1} = aU_{n}+b$. L'énoncé pose ensuite $V_{n} = U_{n}-\alpha$, et
demande de démontrer que $\left(V_{n}\right)$ est géométrique.
\begin{enumerate}[label=\textbf{\arabic*.}, leftmargin=*, itemsep=0.15em]
  \item Écrire $V_{n+1} = U_{n+1}-\alpha$, puis remplacer $U_{n+1}$ par son
        expression.
  \item Faire apparaître $U_{n}-\alpha$ en facteur : on obtient
        $V_{n+1} = a\,V_{n}$.
  \item Conclure : « $\left(V_{n}\right)$ est géométrique de raison $a$ et
        de premier terme $V_{0} = U_{0}-\alpha$ ».
  \item En déduire $V_{n} = V_{0}\,a^{n}$, puis
        $U_{n} = V_{n}+\alpha$.
  \item Pour les sommes : calculer d'abord celle des $V_{n}$, puis ajouter
        $\alpha$ autant de fois qu'il y a de termes.
\end{enumerate}
\end{methode}

\begin{exemple}[le schéma complet, sur le sujet 2026]
Soit $U_{0} = 5$ et $U_{n+1} = \dfrac{U_{n}+3}{4}$ ; on pose
$V_{n} = U_{n}-1$.

\emph{Nature de $\left(V_{n}\right)$.}
\[
  V_{n+1} = U_{n+1}-1 = \frac{U_{n}+3}{4}-1 = \frac{U_{n}-1}{4}
  = \frac14 V_{n} .
\]
La suite $\left(V_{n}\right)$ est géométrique de raison $\frac14$ et de
premier terme $V_{0} = 4$.

\emph{Termes généraux.} $V_{n} = 4\left(\frac14\right)^{n}$, donc
$U_{n} = 1+4\left(\frac14\right)^{n}$.

\emph{Sommes.} D'une part
\[
  S_{n} = V_{0}+\cdots+V_{n}
  = 4\times\frac{1-\left(\frac14\right)^{n+1}}{1-\frac14}
  = \frac{16}{3}\left(1-\left(\frac14\right)^{n+1}\right) ;
\]
d'autre part, comme $U_{k} = V_{k}+1$ et qu'il y a $n+1$ termes,
\[
  T_{n} = U_{0}+\cdots+U_{n} = S_{n}+(n+1) .
\]

\emph{Limite.} $\left(\frac14\right)^{n} \to 0$, donc $U_{n} \to 1$.
\end{exemple}

\begin{propriete}[la variante logarithmique]
Si $\left(U_{n}\right)$ est géométrique à termes strictement positifs, de
raison $q$, alors $W_{n} = \ln\left(U_{n}\right)$ est \textbf{arithmétique}
de raison $\ln q$. Réciproquement, si $\left(U_{n}\right)$ est arithmétique
de raison $r$, alors $V_{n} = e^{U_{n}}$ est géométrique de raison $e^{r}$.
\end{propriete}

\begin{demonstration}
$W_{n+1}-W_{n} = \ln U_{n+1}-\ln U_{n} = \ln\dfrac{U_{n+1}}{U_{n}} = \ln q$,
qui ne dépend pas de $n$.
\end{demonstration}

\begin{entrainement}[suites auxiliaires]
\exo[\niveau{2}] $U_{0} = 4$ et $U_{n+1} = \frac34U_{n}-1$ ; on pose
$V_{n} = U_{n}+4$. Montrer que $\left(V_{n}\right)$ est géométrique.

\exo[\niveau{2}] $U_{0} = -4$ et $U_{n+1} = \frac13U_{n}-4$ ; on pose
$V_{n} = U_{n}+6$. Montrer que $\left(V_{n}\right)$ est géométrique, puis
exprimer $U_{n}$ en fonction de $n$.

\exo[\niveau{3}] $U_{0} = 2$ et $U_{n+1} = 0{,}5U_{n}+3$ ; on pose
$V_{n} = U_{n}-6$. Déterminer $U_{n}$ puis $\limn U_{n}$.

\exo[\niveau{3}] $\left(U_{n}\right)$ arithmétique de premier terme
$U_{0} = 1$ et de raison $-2$ ; on pose $V_{n} = e^{-U_{n}}$. Montrer que
$\left(V_{n}\right)$ est géométrique de raison $e^{2}$.
\end{entrainement}

\begin{corrige}[suites auxiliaires]
\textbf{1.} $V_{n+1} = U_{n+1}+4 = \frac34U_{n}-1+4 = \frac34U_{n}+3
= \frac34\left(U_{n}+4\right) = \frac34V_{n}$ : géométrique de raison
$\frac34$, de premier terme $V_{0} = 8$.
\textbf{2.} $V_{n+1} = \frac13U_{n}-4+6 = \frac13U_{n}+2
= \frac13\left(U_{n}+6\right) = \frac13V_{n}$ : raison $\frac13$,
$V_{0} = 2$. D'où $V_{n} = 2\left(\frac13\right)^{n}$ et
$U_{n} = 2\left(\frac13\right)^{n}-6$.
\textbf{3.} $V_{n+1} = 0{,}5U_{n}+3-6 = 0{,}5\left(U_{n}-6\right)
= 0{,}5V_{n}$ ; $V_{0} = -4$, donc $U_{n} = 6-4\times0{,}5^{n}$, de limite
$6$.
\textbf{4.} $\frac{V_{n+1}}{V_{n}} = e^{-U_{n+1}+U_{n}} = e^{-r} = e^{2}$.
\end{corrige}

% ===========================================================================
\begin{annales}
Ce bloc est, à lui seul, une préparation complète à l'exercice 1 de
l'épreuve. Les énoncés sont repris tels quels.

\exoann{Baccalauréat, Madagascar, série A, session 2014 — exercice 1}
Soit $\left(U_{n}\right)$ définie par $U_{0} = -4$ et
$U_{n+1} = \frac13U_{n}-4$.
\begin{enumerate}[label=\textbf{\arabic*)}, leftmargin=*, itemsep=0.15em]
  \item Calculer $U_{1}$ et $U_{2}$.
  \item Soit $\left(V_{n}\right)$ définie par $V_{n} = U_{n}+6$.
        \begin{enumerate}[label=\textbf{\alph*)}, leftmargin=1.4em]
          \item Démontrer que $\left(V_{n}\right)$ est géométrique, de
                raison $q$ et de premier terme $V_{0}$ à préciser.
          \item Exprimer $V_{n}$ puis $U_{n}$ en fonction de $n$.
        \end{enumerate}
  \item \begin{enumerate}[label=\textbf{\alph*)}, leftmargin=1.4em]
          \item Exprimer en fonction de $n$ les sommes
                $S_{n} = V_{0}+V_{1}+\cdots+V_{n}$ et
                $S'_{n} = U_{0}+U_{1}+\cdots+U_{n}$.
          \item Calculer $\displaystyle\limn \frac{S'_{n}}{n+1}$.
        \end{enumerate}
\end{enumerate}

\exoann{Baccalauréat, Madagascar, série A, session 2026 — exercice 1}
Soit $\left(U_{n}\right)_{n \in \N}$ définie par $U_{0} = 5$ et
$U_{n+1} = \dfrac{U_{n}+3}{4}$.
\begin{enumerate}[label=\textbf{\arabic*)}, leftmargin=*, itemsep=0.15em]
  \item Calculer $U_{1}$ et $U_{2}$.
  \item On pose $V_{n} = U_{n}-1$.
        \begin{enumerate}[label=\textbf{\alph*)}, leftmargin=1.4em]
          \item Montrer que $\left(V_{n}\right)$ est géométrique dont on
                précisera la raison et le premier terme $V_{0}$.
          \item Exprimer $V_{n}$ en fonction de $n$.
          \item En déduire l'expression de $U_{n}$ en fonction de $n$.
        \end{enumerate}
  \item Soit $S_{n} = V_{0}+V_{1}+\cdots+V_{n}$ et
        $T_{n} = U_{0}+U_{1}+\cdots+U_{n}$.
        \begin{enumerate}[label=\textbf{\alph*)}, leftmargin=1.4em]
          \item Exprimer $S_{n}$ en fonction de $n$.
          \item En déduire l'expression de $T_{n}$ en fonction de $n$.
        \end{enumerate}
\end{enumerate}

\exoann{Baccalauréat, Madagascar, série A, session 2021 — exercice 1}
Soit $\left(U_{n}\right)$ définie par $U_{0} = 2$ et, pour tout
$n \in \N$, $2U_{n+1}+1 = U_{n+1}+U_{n}$.
\begin{enumerate}[label=\textbf{\arabic*)}, leftmargin=*, itemsep=0.15em]
  \item \begin{enumerate}[label=\textbf{\alph*)}, leftmargin=1.4em]
          \item Montrer que $\left(U_{n}\right)$ est arithmétique de raison
                $r = -1$.
          \item Exprimer $U_{n}$ en fonction de $n$.
          \item Déterminer $N$ pour que $U_{N} = -98$.
          \item Calculer $S = U_{0}+U_{1}+\cdots+U_{100}$.
        \end{enumerate}
  \item Soit $\left(V_{n}\right)$ définie par $V_{n} = e^{2-n}$. Montrer que
        $\left(V_{n}\right)$ est géométrique dont on déterminera la raison
        $q$ et le premier terme $V_{0}$.
\end{enumerate}

\exoann{Baccalauréat, Madagascar, série A, session 2016 — exercice 1}
Soit $\left(U_{n}\right)$ une suite arithmétique de premier terme
$U_{0} = 1$ et de raison $r = -2$.
\begin{enumerate}[label=\textbf{\arabic*)}, leftmargin=*, itemsep=0.15em]
  \item Exprimer $U_{n}$ en fonction de $n$.
  \item \begin{enumerate}[label=\textbf{\alph*)}, leftmargin=1.4em]
          \item Déterminer l'entier naturel $n$ tel que $U_{n} = -99$.
          \item Soit $S = U_{0}+U_{1}+\cdots+U_{50}$. Prouver que
                $S = -2\,499$.
        \end{enumerate}
  \item Soit $\left(V_{n}\right)$ définie par $V_{n} = e^{-U_{n}}$.
        \begin{enumerate}[label=\textbf{\alph*)}, leftmargin=1.4em]
          \item Démontrer que $\left(V_{n}\right)$ est géométrique de raison
                $q = e^{2}$.
          \item Exprimer $U_{n}$ en fonction de $V_{n}$, puis en déduire
                $\limn U_{n}$.
        \end{enumerate}
\end{enumerate}

\exoann{Baccalauréat, Madagascar, série A, session 2017 — exercice 1}
Pour tout entier naturel $n$, on pose $U_{n} = \left(\frac34\right)^{n}$ et
$V_{n} = \ln\left(U_{n}\right)$.
\begin{enumerate}[label=\textbf{\arabic*)}, leftmargin=*, itemsep=0.15em]
  \item \begin{enumerate}[label=\textbf{\alph*)}, leftmargin=1.4em]
          \item Calculer $U_{0}$, $U_{1}$, $V_{0}$ et $V_{1}$.
          \item Montrer que $\left(U_{n}\right)$ est géométrique de raison
                $q = \frac34$.
          \item Calculer, en fonction de $n$, la somme
                $S_{n} = U_{0}+U_{1}+\cdots+U_{n}$.
        \end{enumerate}
  \item \begin{enumerate}[label=\textbf{\alph*)}, leftmargin=1.4em]
          \item Vérifier que $\left(V_{n}\right)$ est arithmétique, de raison
                à préciser.
          \item En déduire la variation de $\left(V_{n}\right)$.
        \end{enumerate}
  \item Calculer $\limn S_{n}$ et $\limn V_{n}$.
\end{enumerate}

\exoann{Baccalauréat, Madagascar, série A, session 2020 — exercice 1}
\begin{enumerate}[label=\textbf{\arabic*)}, leftmargin=*, itemsep=0.15em]
  \item $\left(V_{n}\right)$ est définie par
        $V_{n} = 3\left(\frac58\right)^{n}$.
        \begin{enumerate}[label=\textbf{\alph*)}, leftmargin=1.4em]
          \item Calculer $V_{0}$ et $V_{1}$.
          \item Exprimer $V_{n+1}$ en fonction de $V_{n}$ ; en déduire la
                nature de $\left(V_{n}\right)$ et préciser sa raison.
          \item Calculer $\limn V_{n}$.
          \item Exprimer en fonction de $n$ la somme
                $S_{n} = V_{0}+V_{1}+\cdots+V_{n-1}$.
        \end{enumerate}
  \item $\left(U_{n}\right)$ est une suite arithmétique telle que
        $U_{34} = U_{15}+57$.
        \begin{enumerate}[label=\textbf{\alph*)}, leftmargin=1.4em]
          \item Vérifier que la raison vaut $3$.
          \item Sachant que $U_{22} = 65$, calculer $U_{34}$ et $U_{15}$.
          \item En déduire $S = U_{15}+U_{16}+\cdots+U_{34}$.
        \end{enumerate}
\end{enumerate}

\exoann{Baccalauréat, Madagascar, série A, session 2015 — exercice 1}
Soit $\left(U_{n}\right)$ une suite arithmétique telle que
$U_{1}+2U_{3}-U_{5} = 2$ et $U_{2} = -1$.
\begin{enumerate}[label=\textbf{\arabic*)}, leftmargin=*, itemsep=0.15em]
  \item \begin{enumerate}[label=\textbf{\alph*)}, leftmargin=1.4em]
          \item Exprimer $U_{1}$, $U_{3}$ et $U_{5}$ en fonction de $U_{2}$
                et de la raison $r$.
          \item Calculer $r$ ; en déduire le sens de variation de
                $\left(U_{n}\right)$.
        \end{enumerate}
  \item Exprimer $U_{n}$ en fonction de $n$ ; calculer $U_{25}$.
  \item Calculer la somme $S = U_{2}+U_{3}+\cdots+U_{25}$.
  \item On pose $V_{n} = e^{3-2n}$ pour tout $n \in \N$.
        \begin{enumerate}[label=\textbf{\alph*)}, leftmargin=1.4em]
          \item Montrer que $\left(V_{n}\right)$ est géométrique dont on
                précisera la raison et le premier terme.
          \item Exprimer la somme $S_{n} = V_{0}+V_{1}+\cdots+V_{n}$ en
                fonction de $n$.
        \end{enumerate}
\end{enumerate}

\exoann{Baccalauréat, Madagascar, série A, session 2013 — exercice 1
(extrait)}
\begin{enumerate}[label=\textbf{\arabic*)}, leftmargin=*, itemsep=0.15em]
  \item $\left(V_{n}\right)$ est géométrique à termes positifs, de premier
        terme $V_{0}$ et de raison $q = \frac34$, telle que
        $V_{3} = \frac{27}{8}$.
        \begin{enumerate}[label=\textbf{\alph*)}, leftmargin=1.4em]
          \item Justifier que $V_{0} = 8$.
          \item Exprimer $V_{n}$ en fonction de $n$.
          \item Déterminer $\limn \left(V_{n}+3\right)$.
        \end{enumerate}
  \item Soit $\left(W_{n}\right)$ arithmétique de terme général
        $W_{n} = \ln\left(8\times\left(\frac34\right)^{n}\right)$ et de
        premier terme $W_{0} = \ln 8$. Déterminer, en fonction de $\ln 2$ et
        de $\ln 3$, la somme $S = W_{0}+W_{1}+\cdots+W_{5}$.
\end{enumerate}
\end{annales}

\begin{corrige}[annales]
\textbf{1.} $U_{1} = -\frac{16}{3}$ et $U_{2} = -\frac{52}{9}$.
$V_{n+1} = \frac13U_{n}-4+6 = \frac13\left(U_{n}+6\right) = \frac13V_{n}$ :
géométrique de raison $\frac13$, $V_{0} = 2$. D'où
$V_{n} = 2\left(\frac13\right)^{n}$ et
$U_{n} = 2\left(\frac13\right)^{n}-6$. Ensuite
$S_{n} = 3\left(1-\left(\frac13\right)^{n+1}\right)$ et
$S'_{n} = S_{n}-6(n+1)$. Enfin
$\frac{S'_{n}}{n+1} = \frac{S_{n}}{n+1}-6 \to -6$, puisque $S_{n}$ tend
vers $3$.

\textbf{2.} Voir l'exemple du cours : $U_{1} = 2$, $U_{2} = \frac54$ ;
$V_{n+1} = \frac14V_{n}$ avec $V_{0} = 4$ ;
$V_{n} = 4\left(\frac14\right)^{n}$ ;
$U_{n} = 1+4\left(\frac14\right)^{n}$ ;
$S_{n} = \frac{16}{3}\left(1-\left(\frac14\right)^{n+1}\right)$ et
$T_{n} = S_{n}+(n+1)$.

\textbf{3.} En retranchant $U_{n+1}$ : $U_{n+1} = U_{n}-1$, d'où
$r = -1$ et $U_{n} = 2-n$. $U_{N} = -98$ donne $N = 100$. La somme compte
$101$ termes : $S = 101\times\frac{2-98}{2} = 101\times(-48) = -4\,848$.
Enfin $\frac{V_{n+1}}{V_{n}} = e^{-1}$ : géométrique de raison $e^{-1}$, de
premier terme $V_{0} = e^{2}$.

\textbf{4.} $U_{n} = 1-2n$ ; $U_{n} = -99$ donne $n = 50$ ; la somme compte
$51$ termes et vaut $51\times\frac{1-99}{2} = -2\,499$. \checkmark
$\frac{V_{n+1}}{V_{n}} = e^{U_{n}-U_{n+1}} = e^{2}$. Comme
$V_{n} = e^{-U_{n}}$, on a $U_{n} = -\ln\left(V_{n}\right)$, et
$U_{n} = 1-2n \to -\infty$.

\textbf{5.} $U_{0} = 1$, $U_{1} = \frac34$, $V_{0} = 0$,
$V_{1} = \ln\frac34 \approx -0{,}29$. La suite $\left(U_{n}\right)$ est
géométrique de raison $\frac34$, donc
$S_{n} = \frac{1-\left(\frac34\right)^{n+1}}{\frac14}
= 4\left(1-\left(\frac34\right)^{n+1}\right)$.
$V_{n+1}-V_{n} = \ln\frac34$, constante : arithmétique de raison
$\ln\frac34<0$, donc décroissante. Enfin $S_{n} \to 4$ et
$V_{n} \to -\infty$.

\textbf{6.} $V_{0} = 3$, $V_{1} = \frac{15}{8}$,
$V_{n+1} = \frac58V_{n}$ : géométrique de raison $\frac58$, limite $0$ ;
$S_{n} = 3\times\frac{1-\left(\frac58\right)^{n}}{\frac38}
= 8\left(1-\left(\frac58\right)^{n}\right)$ (la somme s'arrête à
$V_{n-1}$, donc elle compte $n$ termes). Pour la seconde suite :
$19r = 57$ donc $r = 3$ ; $U_{34} = 101$, $U_{15} = 44$ ; et
$S = 20\times\frac{44+101}{2} = 1\,450$.

\textbf{7.} $U_{1} = U_{2}-r$, $U_{3} = U_{2}+r$, $U_{5} = U_{2}+3r$, donc
$U_{1}+2U_{3}-U_{5} = 2U_{2}-2r = 2$. Avec $U_{2} = -1$ : $-2-2r = 2$, soit
$r = -2$ : la suite est décroissante. Alors
$U_{n} = U_{2}+(n-2)r = -1-2(n-2) = 3-2n$, et $U_{25} = -47$. La somme
compte $24$ termes : $S = 24\times\frac{-1-47}{2} = -576$. Enfin
$\frac{V_{n+1}}{V_{n}} = e^{-2}$ : géométrique de raison $e^{-2}$, de
premier terme $V_{0} = e^{3}$, et
$S_{n} = e^{3}\,\frac{1-e^{-2(n+1)}}{1-e^{-2}}$.

\textbf{8.} $V_{0}q^{3} = V_{0}\times\frac{27}{64} = \frac{27}{8}$ donne
$V_{0} = 8$ ; $V_{n} = 8\left(\frac34\right)^{n}$, de limite $0$, donc
$V_{n}+3 \to 3$. Pour la somme :
$W_{n} = \ln 8+n\ln\frac34$, donc
\[
  S = 6\ln 8+\left(0+1+2+3+4+5\right)\ln\frac34
    = 18\ln 2+15\left(\ln 3-2\ln 2\right) = 15\ln 3-12\ln 2 .
\]
\end{corrige}

% ===========================================================================
\begin{vraifaux}
\exo Une suite géométrique de raison $1$ est aussi arithmétique.

\exo Si $u_{n+1}-u_{n}$ dépend de $n$, la suite n'est pas arithmétique.

\exo La somme des termes d'une suite géométrique de raison $\frac12$ peut
dépasser $10\,u_{0}$.

\exo Si $\left(V_{n}\right)$ est géométrique et $U_{n} = V_{n}+3$, alors
$\left(U_{n}\right)$ est géométrique.

\exo Une suite arithmétique de raison non nulle n'est jamais convergente.
\end{vraifaux}

\begin{corrige}[vrai ou faux]
\textbf{1.} Vrai — elle est constante, donc arithmétique de raison $0$.
\textbf{2.} Vrai — la raison doit être un nombre fixe.
\textbf{3.} Faux — la somme est majorée par $2u_{0}$, puisque
$\frac{1}{1-\frac12} = 2$.
\textbf{4.} Faux — ajouter une constante détruit la structure
géométrique ; c'est précisément pourquoi on pose la suite auxiliaire dans
l'autre sens.
\textbf{5.} Vrai — elle tend vers $+\infty$ ou $-\infty$ selon le signe de
la raison.
\end{corrige}

% ===========================================================================
\begin{situation}[le prêt d'une association villageoise]
Une association villageoise prête $600\,000$ ariary à l'un de ses
membres, au taux mensuel de $2\,\%$. Chaque mois, on ajoute les intérêts au
capital restant dû, puis l'emprunteur verse $50\,000$ ariary. On note
$C_{n}$ le capital restant dû après le $n$-ième versement, avec
$C_{0} = 600\,000$.

\begin{enumerate}[label=\textbf{\arabic*.}, leftmargin=*, itemsep=0.3em]
  \item Justifier que $C_{n+1} = 1{,}02\,C_{n}-50\,000$.
  \item Vérifier que si $C_{n} = 2\,500\,000$, alors
        $C_{n+1} = 2\,500\,000$.
  \item On pose $D_{n} = C_{n}-2\,500\,000$. Montrer que
        $\left(D_{n}\right)$ est géométrique de raison $1{,}02$.
  \item En déduire $C_{n}$ en fonction de $n$.
  \item Calculer $C_{12}$ : où en est l'emprunteur au bout d'un an ?
\end{enumerate}
\end{situation}

\begin{corrige}[le prêt d'une association villageoise]
\textbf{1.} Les intérêts du mois valent $0{,}02\,C_{n}$ ; on les ajoute,
puis on retranche le versement.

\textbf{2.} $1{,}02\times2\,500\,000 = 2\,550\,000$, moins $50\,000$, donne
bien $2\,500\,000$.

\textbf{3.} $D_{n+1} = C_{n+1}-2\,500\,000 = 1{,}02C_{n}-50\,000-2\,500\,000
= 1{,}02\left(C_{n}-2\,500\,000\right) = 1{,}02\,D_{n}$.

\textbf{4.} $D_{0} = -1\,900\,000$, donc
$C_{n} = 2\,500\,000-1\,900\,000\times1{,}02^{n}$.

\textbf{5.} $1{,}02^{12} \approx 1{,}268$, donc
$C_{12} \approx 2\,500\,000-2\,409\,000 = 91\,000$ ariary : au bout d'un an,
il ne reste qu'environ $91\,000$ ariary à rembourser — le prêt sera soldé au
treizième mois.
\end{corrige}

% ===========================================================================
\begin{jemeteste}
Trente minutes.

\begin{enumerate}[label=\textbf{\arabic*.}, leftmargin=*, itemsep=0.3em]
  \item Suite arithmétique de premier terme $5$ et de raison $-3$ : donner
        $u_{n}$, puis $u_{0}+u_{1}+\cdots+u_{20}$.
  \item Suite géométrique de premier terme $2$ et de raison $\frac13$ :
        donner $u_{n}$ et $\limn u_{n}$.
  \item $U_{0} = 7$ et $U_{n+1} = \frac12U_{n}+2$ ; on pose
        $V_{n} = U_{n}-4$. Montrer que $\left(V_{n}\right)$ est géométrique,
        puis donner $U_{n}$ et $\limn U_{n}$.
  \item $V_{n} = e^{4-3n}$ : nature, raison, premier terme, limite.
  \item Combien de termes compte la somme $u_{7}+u_{8}+\cdots+u_{31}$ ?
\end{enumerate}
\end{jemeteste}

\begin{corrige}[je me teste]
\textbf{1.} $u_{n} = 5-3n$ ; $21$ termes, dernier $u_{20} = -55$ :
$S = 21\times\frac{5-55}{2} = -525$.
\textbf{2.} $u_{n} = 2\left(\frac13\right)^{n}$, de limite $0$.
\textbf{3.} $V_{n+1} = \frac12U_{n}+2-4 = \frac12\left(U_{n}-4\right)
= \frac12V_{n}$ ; $V_{0} = 3$, donc
$U_{n} = 4+3\left(\frac12\right)^{n} \to 4$.
\textbf{4.} $\frac{V_{n+1}}{V_{n}} = e^{-3}$ : géométrique de raison
$e^{-3}$, premier terme $e^{4}$, limite $0$.
\textbf{5.} $31-7+1 = 25$ termes.
\end{corrige}

% ===========================================================================
\begin{commeaubac}
\bmtsource{D'après le baccalauréat, Madagascar, série A, session 2022 —
exercice 1 \textbullet\ 5 points}
\emph{La copie de ce sujet dont on dispose est partiellement illisible :
l'énoncé a été rétabli à partir de ce qui s'y lit, et chacun des résultats
annoncés y est redémontré ci-dessous.}

On considère la suite numérique $\left(U_{n}\right)_{n \in \N}$ définie par
\[
  U_{0} = -1, \qquad U_{n+1} = \frac14 U_{n}+3 \quad
  \text{pour tout } n \in \N .
\]

\begin{enumerate}[label=\textbf{\arabic*.}, leftmargin=*, itemsep=0.25em]
  \item Calculer $U_{1}$ et $U_{2}$. \bareme{0,25 $\times$ 2 pts}
  \item Soit la suite numérique $\left(V_{n}\right)_{n \in \N}$ définie par
        $V_{n} = U_{n}-4$ pour tout $n \in \N$.
    \begin{enumerate}[label=\textbf{\alph*)}, leftmargin=1.4em, itemsep=0.15em]
      \item Montrer que $\left(V_{n}\right)$ est une suite géométrique dont
            on précisera la raison et le premier terme. \bareme{1,5 pt}
      \item Exprimer $V_{n}$ puis $U_{n}$ en fonction de $n$.
            \bareme{0,5 pt $\times$ 2}
    \end{enumerate}
  \item On donne $W_{n} = \ln\left(5\left(\frac14\right)^{n}\right)$ pour
        tout $n \in \N$.
    \begin{enumerate}[label=\textbf{\alph*)}, leftmargin=1.4em, itemsep=0.15em]
      \item Montrer que $\left(W_{n}\right)$ est une suite arithmétique dont
            on précisera la raison et le premier terme. \bareme{1,5 pt}
      \item Calculer $\displaystyle\limn W_{n}$. \bareme{0,5 pt}
    \end{enumerate}
\end{enumerate}
\end{commeaubac}

\begin{corrige}[comme au baccalauréat — Madagascar A 2022]
\textbf{1.} $U_{1} = -\frac14+3 = \frac{11}{4}$ et
$U_{2} = \frac{11}{16}+3 = \frac{59}{16}$.

\textbf{2. a)}
\[
  V_{n+1} = U_{n+1}-4 = \frac14U_{n}+3-4 = \frac14U_{n}-1
  = \frac14\left(U_{n}-4\right) = \frac14 V_{n} .
\]
La suite $\left(V_{n}\right)$ est donc géométrique de raison $q = \frac14$
et de premier terme $V_{0} = U_{0}-4 = -5$.
\textbf{b)} $V_{n} = -5\left(\frac14\right)^{n}$, d'où
$U_{n} = 4-5\left(\frac14\right)^{n}$. On vérifie : $U_{0} = 4-5 = -1$
\checkmark et $U_{1} = 4-\frac54 = \frac{11}{4}$ \checkmark.

\textbf{3. a)} En utilisant les propriétés du logarithme,
\[
  W_{n} = \ln 5+n\ln\frac14,
\]
donc $W_{n+1}-W_{n} = \ln\frac14$, qui ne dépend pas de $n$ : la suite
$\left(W_{n}\right)$ est arithmétique de raison $r = \ln\frac14 = -2\ln2$ et
de premier terme $W_{0} = \ln 5$.
\textbf{b)} La raison étant strictement négative,
$\limn W_{n} = -\infty$.
\end{corrige}

\begin{notepourlaclasse}
Quatre heures ici, deux au chapitre précédent : c'est la répartition qui
donne les meilleurs résultats. Tout se joue sur la méthode de la suite
auxiliaire, qui doit être rédigée au moins six fois dans l'année, toujours
dans le même ordre — trois lignes de calcul, une phrase de conclusion avec
la raison \emph{et} le premier terme.

Le comptage des termes d'une somme mérite une séance à lui seul. Faites
écrire au tableau les trois cas — de $u_{0}$ à $u_{n}$, de $u_{1}$ à
$u_{n}$, de $u_{p}$ à $u_{n}$ — et faites-les recopier dans le cahier de
formules. C'est là que se perdent, chaque année, les demi-points les plus
faciles à sauver.

Le bloc d'annales contient huit sujets réels : un par séance de révision au
troisième trimestre, avec correction en classe des deux plus récents. Un
élève qui les aura tous vus aura, de fait, révisé l'exercice 1 de son propre
examen.
\end{notepourlaclasse}
